%% toc: P15 163 163.10
% Comparison of the exchangeability route with the inductive task (Zabell HHL 10, pp. 272--283, 289;
% Durrett Version 5, Section 4.7) and with H\'ajek--Staffel, Handbook of Rationality ch. 4.1 (pp. 239--248).
\chapter{Exchangeability Assumptions and Their Limits as a Justification of Induction}
\label{sec:p15-w02b-p15-022}

This chapter asks what the de Finetti representation contributes to a
particular inductive task: predicting further instances of a kind from
observed instances. The mathematical results it uses are proved in
\cref{sec:p15-exchangeability-and-definetti}; the philosophical setting of
subjective probability is that of \cref{sec:p15-rationality-4-1}. Here the
two are compared directly. The chapter states (i) what supports adopting an
exchangeable law for the task, (ii) what the symmetry and mixture
conclusions then deliver for that task, and (iii) which failures and
objections limit that delivery as a justification of induction.

The \emph{inputs} are a description of the instances and their outcomes,
an order in which instances are indexed, a joint probability law over all
finite records, and the observed initial record. The \emph{conclusions}
are predictive probabilities for later instances and statements about
limiting relative frequencies. Every conclusion below is conditional on the
stated law. None of the theorems makes a physical sequence exchangeable,
and none supplies the normative authority of updating by conditioning
\cite[pp.~249--255]{DurrettPTEAuthor2019}.

\section{The inductive task and its inputs}
\label{sec:p15-w02b-p15-022-task}

\begin{hbdefinition}{def:p15-w02b-p15-022-task}{An exchangeable binary prediction task}
Instances are indexed \(1,2,\ldots\) before any outcome is observed. The
outcome of instance \(j\) is \(X_j\in\{0,1\}\), where \(1\) records that
the instance has the property under study. Work on the canonical space
\(\Omega=\{0,1\}^{\mathbb N_+}\) with its product sigma-field \(\mathcal F\),
coordinates \(X_j\), and a probability law \(P\). Put
\(S_n=\sum_{j=1}^nX_j\), let \(\mathcal F_n=\sigma(X_1,\ldots,X_n)\),
with \(\mathcal F_0\) trivial, and write \(X_{1:n}=b\) for the event that
the first \(n\) outcomes form the word \(b\in\{0,1\}^n\).

The \emph{exchangeability assumption} is that \(P\) is exchangeable in the
sense of \cref{def:p15-exchangeable-canonical}, with exchangeable
sigma-field \(\mathcal E\). Its \emph{directing frequency} \(\Theta\) is a
version of \(E(X_1\mid\mathcal E)\) with values in \([0,1]\), and its
\emph{mixing law} \(\lambda\) is the law of \(\Theta\). By
\cref{thm:p15-binary-definetti,thm:p15-exchangeable-strong-laws},
\(\lambda\) is the unique mixing measure of the binary representation and
\(S_n/n\to\Theta\) almost surely. For a probability \(\lambda\) on
\([0,1]\), \(P_\lambda\) denotes the exchangeable law with that mixing law.

The \emph{task} is to state, from an observed word \(X_{1:n}=b\), a
probability for \(X_{N+1}=1\) with \(N\geq n\), or a claim about the
limiting frequency. The binary case gives explicit predictive formulas;
the general standard-space mixture is \cref{thm:p15-definetti-kernel}.
\index{exchangeability!inductive task}\index{directing frequency}
\end{hbdefinition}

The order of indexing is part of the input. An exchangeable law treats
that order as carrying no information about the outcomes; it does not say
that the order of the actual world carries none.

\section{What supports adopting the symmetry}
\label{sec:p15-w02b-p15-022-support}

Zabell locates the inductive content of the permutation postulate in a
comparison with sampling. Observed instances are like draws from an urn
only to the degree that they are judged exchangeable. Repetition of
instances is not enough; the critics of the rule of succession argued that
some further assumption is needed, and exchangeability is a precise form of
that further assumption. Zabell lists Mill's uniformity of nature, Keynes's
limited variety and Goodman's projectibility as earlier names for such a
principle \cite[p.~272]{Zabell2011CarnapInductiveInference}.

What could support the judgment is therefore a reason to think that the
labels \(1,2,\ldots\) carry no information about the outcomes: for example
that the procedure assigning indices to instances is the same for every
instance, and that nothing distinguishing earlier from later instances is
relevant to the property recorded. Such a reason concerns the particular
task and its description. It is not supplied by probability theory. The
following consequences show what the judgment commits one to, and so what
evidence can tell against it.

\begin{hblemma}{thm:p15-w02b-p15-022-symmetry-consequences}{proved=annex:p15-w02b-p15-022-symmetry-consequences}{What exchangeability commits a law to}
Let \((X_1,\ldots,X_n)\), \(n\geq2\), be an exchangeable vector with values
in a measurable space \(S\).
\begin{enumerate}
\item Every \(X_j\) has the law of \(X_1\).
\item For \(i\neq j\), the pair \((X_i,X_j)\) has the law of \((X_1,X_2)\),
      and \((X_2,X_1)\) has that same law.
\item For an infinite exchangeable binary sequence and every \(k\geq1\),
      \(P(X_1=\cdots=X_k=1)=E\Theta^k\). In particular
      \(P(X_1=1)=E\Theta\) and \(P(X_1=X_2=1)=E\Theta^2\).
\end{enumerate}
\end{hblemma}

Item~1 excludes any trend in the probability of the property along the
index. Item~2 excludes any dependence that varies with the distance
between instances. Neither consequence is implied by repeated observation of
the same kind of instance. Item~3 makes the moments of the mixing law
observable as probabilities of finite records.

\section{What the symmetry and the mixture deliver}
\label{sec:p15-w02b-p15-022-conclusion}

\begin{hbtheorem}{thm:p15-w02b-p15-022-predictive}{proved=annex:p15-w02b-p15-022-predictive}{Predictive updating and limiting frequency under exchangeability}
Let \(P\) be an infinite exchangeable binary law with directing frequency
\(\Theta\) and mixing law \(\lambda\).
\begin{enumerate}
\item For integers \(N\geq n\geq0\),
      \(E(X_{N+1}\mid\sigma(\mathcal F_n\cup\mathcal E))=\Theta\) almost
      surely. Consequently
      \[
        P(X_{N+1}=1\mid\mathcal F_n)=E(\Theta\mid\mathcal F_n)
                                         \quad\text{almost surely},
      \]
      and the same version serves every later index \(N+1\).
\item If \(b\in\{0,1\}^n\) has \(s\) ones and \(P(X_{1:n}=b)>0\), then
      \[
       P(X_{n+1}=1\mid X_{1:n}=b)
          =\frac{\int_0^1\theta^{s+1}(1-\theta)^{n-s}\,\lambda(d\theta)}
                {\int_0^1\theta^{s}(1-\theta)^{n-s}\,\lambda(d\theta)}.
      \]
      It depends on \(b\) only through \(n\) and \(s\).
\item Put \(Q_n=P(X_{n+1}=1\mid\mathcal F_n)\). Then \(Q_n\to\Theta\)
      almost surely and in \(L^1\), and \(S_n/n\to\Theta\) almost surely.
      Hence \(Q_n-S_n/n\to0\) almost surely.
\end{enumerate}
\end{hbtheorem}

This is what the mixture conclusion contributes to the task. Given the
exchangeable field, the outcomes behave as independent trials with success
probability \(\Theta\); the agent's predictive probability is the
conditional mean of the unknown \(\Theta\); and predictive probabilities and
observed frequencies approach the same limit. Zabell's dependence on counts
alone is item~2 \cite[p.~273]{Zabell2011CarnapInductiveInference}.
Durrett calls symmetry a palatable assumption with a powerful conclusion
\cite[p.~255]{DurrettPTEAuthor2019}. He offers the binary representation
as useful to anyone concerned with the foundations of statistics, and for
that use he refers to Section~3.7 of the second edition of Savage's
\emph{The Foundations of Statistics} (Dover, 1972)
\cite[p.~255]{DurrettPTEAuthor2019}. Savage's discussion is not held for
this handbook and is not reconstructed here; the remark is Durrett's
report of a use, not a theorem. Item~3, however, is an almost-sure
statement under \(P\) itself. It says that the agent is certain, by the
agent's own law, that predictions will track frequencies. It says nothing
about a record that the law treats as negligible.

\begin{hbtheorem}{thm:p15-w02b-p15-022-mixing-law}{proved=annex:p15-w02b-p15-022-mixing-law}{The same symmetry with different mixing laws}
Fix \(x\in\{0,1\}^{\mathbb N_+}\) whose initial counts \(s_n\) satisfy
\(s_n/n\to f\in[0,1]\). Write \(p_n(\lambda)\) for
\(P_\lambda(X_{n+1}=1\mid X_{1:n}=x_{1:n})\) whenever the conditioning word
has positive probability.
\begin{enumerate}
\item If \(\lambda=\delta_p\), then \(p_n(\lambda)=p\) whenever it is defined.
\item If \(\lambda\) is uniform on \([0,1]\), every word has positive
      probability and \(p_n(\lambda)=(s_n+1)/(n+2)\to f\). More generally,
      for \(\lambda=\operatorname{beta}(\alpha,\beta)\),
      \(p_n(\lambda)=(\alpha+s_n)/(\alpha+\beta+n)\).
\item If \(\lambda=\frac12\delta_{1/10}+\frac12\delta_{9/10}\), every word has
      positive probability, \(p_n(\lambda)\to9/10\) when \(f>1/2\), and
      \(p_n(\lambda)\to1/10\) when \(f<1/2\).
\item Under \(P_\lambda\), \(\Theta\) lies in the support of \(\lambda\)
      almost surely. Hence the set of sequences whose frequency converges
      to a value outside that support is \(P_\lambda\)-null.
\end{enumerate}
\end{hbtheorem}

All four laws are exchangeable. With \(f=7/10\), the uniform law predicts
asymptotically \(7/10\), the two-point law \(9/10\), and the point mass at
\(1/2\) always \(1/2\). Exchangeability therefore fixes the form of
learning, not how strongly or toward what one learns. Zabell stresses that
the representation places no restriction on the mixing measure
\cite[pp.~275--276]{Zabell2011CarnapInductiveInference}. The point mass is
the independent case, in which the past carries no information about the
future \cite[p.~282]{Zabell2011CarnapInductiveInference}.

A Johnson--Carnap predictive rule, or equivalently a Dirichlet mixing law,
needs content beyond symmetry. \Cref{thm:p15-carnap-history-family} gives an
exchangeable and category-symmetric die law whose predictions depend on
counts other than the predicted category's own; Zabell's
\((20,10,10,10,10,10)\) and \((20,30,5,5,5,5)\) comparison makes the same
point \cite[pp.~280--281]{Zabell2011CarnapInductiveInference}.

\section{Failures and objections that limit the justification}
\label{sec:p15-w02b-p15-022-limits}

\begin{hbwarning}{def:p15-w02b-p15-022-not-established}{The representation does not establish the symmetry}
Each theorem above assumes that \(P\) is exchangeable. None infers
exchangeability from observed instances, and none shows that a physical
process has an exchangeable description. Hume's demand for a connexion
between observed and unobserved instances is therefore relocated into the
symmetry judgment, not answered by the representation; see
\cref{sec:p15-hume-reference-use}. A finite record cannot settle the
judgment either, as the next theorem shows.
\end{hbwarning}

\begin{hbtheorem}{thm:p15-w02b-p15-022-recoding}{proved=annex:p15-w02b-p15-022-recoding}{Exchangeability depends on the description of the outcomes}
Let \(X\) be an infinite exchangeable binary sequence under \(P\), let
\(N\geq1\), and define
\[
 Y_j=X_j\quad(j\leq N),\qquad Y_j=1-X_j\quad(j>N).
\]
\begin{enumerate}
\item The records \(X_{1:N}\) and \(Y_{1:N}\) coincide, so both descriptions
      assign the same probability to every observation of the first \(N\)
      instances.
\item \(Y\) is exchangeable under \(P\) if and only if the \(X_j\) are
      independent with \(P(X_j=1)=1/2\).
\item Let \(P'\) be a second law under which \(Y\) is exchangeable with
      uniform mixing law. If every \(X_j\), \(j\leq N\), equals one, then
      \[
       \begin{aligned}
       P_{\mathrm{unif}}(X_{N+1}=1\mid X_{1:N}=1\cdots1)&=\frac{N+1}{N+2},\\
       P'(X_{N+1}=1\mid X_{1:N}=1\cdots1)&=\frac1{N+2},
       \end{aligned}
      \]
      where \(P_{\mathrm{unif}}\) makes \(X\) exchangeable with uniform
      mixing law. The two laws agree on every event determined by the first
      \(N\) instances.
\end{enumerate}
\end{hbtheorem}

With \(X_j\) recording green and \(Y_j\) recording grue for an emerald
inspected in order \(j\), and every emerald assumed either green or blue,
\(Y\) is the fixed-cutoff predicate of \cref{def:p15-goodman-grue}. Item~2
shows that one law cannot treat both descriptions as exchangeable unless it
is the fair independent law, under which nothing is learned. Item~3 shows
that two agents, each adopting exchangeability for one description and the
same uniform mixing law, agree on all \(N\) records and then make opposite
predictions. The judgment of which description is exchangeable is the
projectibility judgment of \cref{sec:p15-goodman-grue}, not a consequence of
the representation.

\subsection{Finite and nonextendable sequences}

For a finite population sampled without replacement, the exchangeable law
is a mixture of hypergeometric sampling laws, not of independent trials.
\Cref{thm:p15-carnap-history-finite} proves this finite representation,
gives an exchangeable two-observation law with no exchangeable extension to
three, and shows that only extendability to arbitrary length forces
nonnegative instantial relevance. In the three-red, two-black urn, a first
red lowers the probability of a second red from \(3/5\) to \(1/2\).
\Cref{thm:p15-exchangeable-second-moments} gives the corresponding negative
covariance. Zabell states the equivalence that a finite exchangeable
sequence is infinitely extendable exactly when it is a mixture of
independent identically distributed sequences, and he insists that
nonextendable finite populations are legitimate cases, not pathologies
\cite[pp.~281--283]{Zabell2011CarnapInductiveInference}. For such a task the
conclusions of \cref{thm:p15-w02b-p15-022-predictive} are unavailable;
what remains is the finite mixture and its own predictive rule.

\subsection{Order, trend and Markov dependence}

\begin{hbtheorem}{thm:p15-w02b-p15-022-markov-trend}{proved=annex:p15-w02b-p15-022-markov-trend}{A stationary chain and a trend violate the symmetry}
For \(q\in[0,1]\), the finite laws
\[
 P_q(X_{1:n}=x)=\frac12\prod_{j=1}^{n-1}
   \bigl(q\mathbf1_{\{x_{j+1}=x_j\}}+(1-q)\mathbf1_{\{x_{j+1}\neq x_j\}}\bigr)
\]
are consistent and define a law on \(\{0,1\}^{\mathbb N_+}\) with
\(P_q(X_j=1)=1/2\) for every \(j\).
\begin{enumerate}
\item \(P_q\) is exchangeable if and only if \(q\in\{1/2,1\}\).
\item Any law with \(P(X_i=1)\neq P(X_j=1)\) for some \(i\neq j\) is not
      exchangeable.
\item For \(0<q<1\), \(P_q(X_5=1\mid X_{1:4}=0011)=q\) and
      \(P_q(X_5=1\mid X_{1:4}=1100)=1-q\). Every law under which
      \((X_1,\ldots,X_5)\) is exchangeable and both words have positive
      probability gives the two words the same predictive probability.
\end{enumerate}
\end{hbtheorem}

Identical one-dimensional laws, as in this stationary chain, therefore do
not suffice for exchangeability. When successive instances persist,
the record \(0011\) and the record \(1100\) carry different information
about the next instance; an exchangeable law discards that information by
construction. With \(q=9/10\) the chain predicts \(9/10\) and \(1/10\),
whereas the uniform exchangeable law predicts \(1/2\) for both records.
Zabell reports de Finetti's partial exchangeability as the response:
Markov exchangeability requires equal probability only for records with the
same initial state and transition counts, and recurrent Markov exchangeable
sequences are mixtures of Markov chains, a result he attributes to Diaconis
and Freedman. He also reports their remark that the appropriate symmetry is
recognized only after experience and reflection
\cite[pp.~289--290]{Zabell2011CarnapInductiveInference}. That representation
is reported here, not proved.

\subsection{Dependence on the mixing law and on the representation}

\Cref{thm:p15-w02b-p15-022-mixing-law} shows that one symmetry is compatible
with no learning, with convergence to every frequency, and with convergence
to a wrong value. \Cref{thm:p15-w02b-p15-022-recoding} shows that the
symmetry itself depends on how the outcomes are described. Both limitations
recur in the comparison below, where they meet the representation
dependence that H\'ajek and Staffel raise against the Indifference
Principle.

\section{Comparison with subjective probability and its dynamics}
\label{sec:p15-w02b-p15-022-comparison}

H\'ajek and Staffel survey probabilism, its betting, representation and
accuracy arguments, Regularity, the Principal Principle, Reflection,
Indifference, conditioning, Jeffrey conditioning and minimal-change rules
\cite[pp.~239--248]{HajekStaffel2021Dynamics}. Their chapter does not
discuss exchangeability or the de Finetti representation theorem; de Finetti
appears there for the betting account, the Dutch-book theorem and an
accuracy-dominance result. The comparison in this section is therefore an application of
their stated norms and objections to the exchangeable task, made here. The
reconstruction of their chapter itself is \cref{sec:p15-rationality-4-1}.

\subsection{Coherence and the source's finite additivity}

\begin{hbtheorem}{thm:p15-w02b-p15-022-finite-word-credences}{proved=annex:p15-w02b-p15-022-finite-word-credences}{What coherent finite-record credences fix}
Let \(\mathcal A\) be the algebra of events \(\{X_{1:n}\in B\}\),
\(n\geq1\), \(B\subseteq\{0,1\}^n\), and let \(C\) be a finitely additive
probability on \(\mathcal A\).
\begin{enumerate}
\item There is exactly one countably additive probability \(P\) on
      \(\mathcal F\) that agrees with \(C\) on \(\mathcal A\). It is
      exchangeable if and only if, for every \(n\), \(C\) gives each word of
      length \(n\) the same value as each of its rearrangements. In that
      case the predictive formula of
      \cref{thm:p15-w02b-p15-022-predictive} holds for \(C\) on
      \(\mathcal A\).
\item The event \(L\) that \(S_n/n\) converges, and its complement, contain
      no nonempty member of \(\mathcal A\). Hence approximation of \(L\)
      from inside and from outside by members of \(\mathcal A\) gives only
      the bounds \(0\) and \(1\), while \(P(L)=1\) for every exchangeable
      \(P\).
\item Every \(P_q\) of \cref{thm:p15-w02b-p15-022-markov-trend} restricts
      to a finitely additive probability on \(\mathcal A\); for
      \(q\notin\{1/2,1\}\) it is not exchangeable.
\end{enumerate}
\end{hbtheorem}

\paragraph{Betting and Dutch-book coherence.}
The Dutch-book theorem and its converse concern whether betting prices obey
the probability calculus \cite[pp.~240--241]{HajekStaffel2021Dynamics}; the
finite equivalence is \cref{thm:p17-finite-price-coherence}. By item~3,
coherent credences need not be exchangeable, and by the converse theorem a
nonexchangeable probability is as immune to a Dutch book as an exchangeable
one. Coherence arguments, betting or accuracy based, therefore support
representing the joint credence by some probability. They do not select an
exchangeable one, and they do not select a mixing law. The objections H\'ajek
and Staffel record against the betting identification of credence apply
unchanged to credences about records.

\paragraph{Finite additivity.}
H\'ajek and Staffel do not assume countable additivity
\cite[p.~248, n.~3]{HajekStaffel2021Dynamics}. Item~1 shows that, for
credences about finite records, nothing is lost: the predictive formula is
a statement about finite-record events. The limiting-frequency conclusions
of \cref{thm:p15-w02b-p15-022-predictive}, item~3, are different. Item~2
shows that they concern an event that finite records neither contain nor
exclude, and the value one comes from the countably additive extension. An
agent who accepts only finite additivity has not thereby accepted those
conclusions.

\paragraph{Conditioning as a proposed norm.}
The predictive formula is a ratio of the agent's present credences. That the
agent's later credence ought to equal it is conditionalization, a proposed
diachronic norm whose support H\'ajek and Staffel present as contested:
diachronic Dutch-book arguments with their converse, objections from the
comparison with disagreement between different people and from the
possibility of declining bets, and expected-accuracy arguments that may
support only planning to conditionalize \cite[p.~244]{HajekStaffel2021Dynamics}.
See \cref{sec:p15-rat41-dynamics}.

\begin{hbexample}{def:p15-w02b-p15-022-conditioning-example}{The representation does not compel the update}
Take the uniform mixing law and the record \(X_{1:3}=111\). The present
credence that \(X_4=1\) conditional on that record is \(4/5\), by
\cref{thm:p15-w02b-p15-022-mixing-law}. Suppose that after observing it the
agent adopts the independent fair law for \(X_4,X_5,\ldots\), so that the
new credence in \(X_4=1\) is \(1/2\). Each credence function, taken at its
own time, is a probability and hence synchronically coherent. Only the pair
violates conditioning. The representation of the earlier law is unaffected;
it described what conditioning would yield. Whether the agent was required
to yield it is the question of the diachronic arguments above.
\end{hbexample}

\paragraph{Indifference and representation dependence.}
Exchangeability is a symmetry among instances, not among outcomes or
frequencies. It does not imply the uniform mixing law of Bayes and Laplace;
adding that law is a further indifference judgment, here over the unknown
frequency. H\'ajek and Staffel present the factory producing cubes of side
at most one foot: uniformity in side, in face area and in volume give
\(1/2\), \(1/4\) and \(1/8\) for a side of at most half a foot, and no
description is privileged \cite[pp.~243--244]{HajekStaffel2021Dynamics}.
The same structure appears in the mixing law.

\begin{hbexample}{def:p15-w02b-p15-022-indifference-example}{Two coordinates of indifference}
Uniformity for \(\Theta\) gives \(\lambda\) with density \(1\); uniformity
for \(\eta=\Theta^2\) gives \(\Theta\) the density \(2\theta\) on
\((0,1)\), by \cref{thm:p15-bayes-transformation}. Both mixing laws define
exchangeable sequences. By \cref{thm:p15-w02b-p15-022-predictive}, item~2,
the first-instance probabilities are
\[
 \int_0^1\theta\,d\theta=\frac12,\qquad
 \int_0^1\theta\cdot2\theta\,d\theta=\frac23,
\]
and after one observed success they are
\[
 \frac{\int_0^1\theta^2\,d\theta}{\int_0^1\theta\,d\theta}=\frac23,\qquad
 \frac{\int_0^1\theta^2\cdot2\theta\,d\theta}{\int_0^1\theta\cdot2\theta\,d\theta}
   =\frac{1/2}{2/3}=\frac34 .
\]
These agree with the posterior means in
\cref{thm:p15-bayes-transformation}. The symmetry among instances is
common to both; the disagreement lies wholly in which coordinate of the
unknown frequency is treated indifferently.
\end{hbexample}

\Cref{thm:p15-w02b-p15-022-recoding} is a second, sharper form of
representation dependence: there the symmetry among instances itself, not
only a prior over frequencies, depends on the description.

\paragraph{Reflection.}
Reflection asks present credence to defer to anticipated future credence,
and H\'ajek and Staffel record that the principle must be restricted when
the future self is expected to be irrational or forgetful
\cite[p.~243]{HajekStaffel2021Dynamics}.

\begin{hbtheorem}{thm:p15-w02b-p15-022-reflection}{proved=annex:p15-w02b-p15-022-reflection}{Deference to a future conditioning self, and its failure}
Let \(P\) be infinite exchangeable binary. For \(N\geq m\geq0\), let
\(Q^N_m\) be the version of \(P(X_{N+1}=1\mid\mathcal F_m)\) equal to the
ratio on each positive-probability event \(\{X_{1:m}=b\}\) and to zero on
each null one.
\begin{enumerate}
\item For \(0\leq n\leq m\leq N\),
      \(E(Q^N_m\mid\mathcal F_n)=Q^N_n\) and \(Q^N_m=Q^m_m\) almost surely.
\item If \(n\leq m\), \(c\in\{0,1\}^n\), and
      \(P(X_{1:n}=c,\ Q^N_m=x)>0\), then
      \[
             P(X_{N+1}=1\mid X_{1:n}=c,\ Q^N_m=x)=x .
      \]
\item Let \(P\) be the independent fair law and let the future credence
      be \(R_1=P'(X_2=1\mid X_1)\) for the uniform exchangeable law \(P'\).
      Then \(R_1=2/3\) exactly on \(\{X_1=1\}\), but
      \(P(X_2=1\mid R_1=2/3)=1/2\).
\end{enumerate}
\end{hbtheorem}

An agent with an exchangeable law who is certain to conditionalize that same
law, and to forget nothing, satisfies Reflection toward that future self
automatically: item~2 is the Reflection identity, with the present credence
conditioned on any present record. The predictive probabilities form a
martingale. The identity needs the agent to know the future updating rule,
so that the proposition about future credence becomes an event about
records. Item~3 shows the restriction H\'ajek and Staffel require: deference
fails toward a future self who will use a different mixing law, which the
present self regards as miscalibrated. Reflection thus links the present
mixing law with the anticipated one; it does not certify either.

\paragraph{Chance.}
Item~1 of \cref{thm:p15-w02b-p15-022-predictive} has the shape of the
Principal Principle: given the directing frequency and past outcomes,
credence in a further success equals \(\Theta\). It is an identity of
conditional expectations, not of ratio conditionals, because
\(\{\Theta=x\}\) typically has probability zero. Treating \(\Theta\) as a
chance is a further interpretive step; H\'ajek and Staffel mention frequency,
propensity and best-system accounts of chance without adopting one
\cite[pp.~240 and~243]{HajekStaffel2021Dynamics}.

\begin{table}[htbp]
\centering\small
\begin{tabular}{p{0.2\linewidth}p{0.36\linewidth}p{0.36\linewidth}}
\hline
Surveyed norm & Bearing on adopting an exchangeable law & Bearing on the mixture and predictive conclusions\\
\hline
Probabilism (betting, accuracy) & Requires some probability; nonexchangeable laws are equally coherent & None beyond the existence of the law that is represented\\
Finite additivity only & Finite-record credences still determine an exchangeable extension & Predictive formula survives; limiting-frequency claims need the countably additive extension\\
Conditioning & No bearing on the symmetry & Needed to turn the predictive ratio into later credence; its justification is contested\\
Indifference & Distinct: symmetry of instances, not of outcomes; also description dependent & A uniform mixing law is an extra indifference judgment, coordinate dependent\\
Reflection & No bearing on the symmetry & Holds toward a future self conditioning the same law; fails toward a differently mixing self\\
Principal Principle & No bearing on the symmetry & Formal analogue with \(\Theta\) as chance, if that identification is accepted\\
\hline
\end{tabular}
\caption{How the surveyed norms bear on the exchangeable inductive task.}
\label{tab:p15-w02b-p15-022-comparison}
\end{table}

\section{Alternatives to probabilistic belief change}
\label{sec:p15-w02b-p15-022-alternatives}

H\'ajek and Staffel close their survey by naming, among the topics they did
not treat, the comparison of Bayesianism with Dempster--Shafer functions,
ranking functions, belief revision theory and possibility theory
\cite[p.~248]{HajekStaffel2021Dynamics}. Their note~5 directs a further
putative argument for probabilism, Cox's theorem, to Dubois and Prade's
chapter of the same handbook \cite[p.~248, n.~5]{HajekStaffel2021Dynamics}.
The source gives none of these frameworks a definition, an update rule or a
comparison argument, and its Table~4.1.1 has rows only for probabilistic
norms. The comparison below is therefore made here, under the source's own
assumptions: precise real-valued credences, finite additivity, conditioning
on evidence learned with certainty and of positive prior probability, and
Jeffrey updating on a partition whose cells have positive old weight
(\cref{def:p15-rat41-jeffrey}).

\begin{longtable}{@{}p{.16\linewidth}p{.24\linewidth}p{.24\linewidth}p{.24\linewidth}@{}}
\caption{Where the alternatives named by H\'ajek and Staffel are treated.}
\label{tab:p15-w02b-p15-022-alternatives}\\
\toprule Framework & Handbook treatment & Static object and certain learning & What the survey's assumptions exclude \\
\midrule\endfirsthead
\caption[]{Where the alternatives named by H\'ajek and Staffel are treated (continued)}\\
\toprule Framework & Handbook treatment & Static object and certain learning & What the survey's assumptions exclude \\
\midrule\endhead
\midrule\multicolumn{4}{r}{\emph{continued on next page}}\\
\endfoot
\endlastfoot
Dempster--Shafer functions &
\cref{sec:p15-r47-evidence}: \cref{def:p15-r47-mass,def:p15-r47-combination},
\cref{thm:p15-r47-focal,thm:p15-r47-combination,thm:p15-r47-pignistic} &
Basic mass, with nonadditive \(\operatorname{Bel}\) and
\(\operatorname{Pl}\); learning \(E\) as combination with mass one on
\(E\); uncertain testimony as simple support &
Additivity; a unique two-sided betting price; the conditioning of each
member of a probability set \\
Ranking functions &
\cref{sec:p15-rationality-5-3}: \cref{def:p15-rat53-ranking,def:p15-rat53-conditionals},
\cref{def:p15-rat53-soft-update,thm:p15-rat53-update}; the bridge
\cref{def:p15-r47-ranking,thm:p15-r47-ranking} &
Integer disbelief ranks; learning \(E\) as \(\kappa(\cdot\mid E)\); learning
with finite firmness by preserving within-cell conditional ranks &
Real-valued additive degrees; the betting and accuracy routes of the
survey's table \\
Belief revision &
\cref{sec:p08-paraconsistent-and-nonmonotonic}:
\cref{def:p08-agm,def:p08-agm-faithful-preorder},
\cref{thm:p08-agm-representation} &
Deductively closed belief set; revision \(K*A\) by minimal-world
selection; no graded input &
Graded credence; iterated revision is outside the eight postulates \\
Possibility theory &
\cref{sec:p15-r47-modal}, \cref{thm:p15-r47-finite-modal},
\cref{thm:p15-r47-ranking}, \cref{thm:p15-ip-possibility-bridge} &
Maxitive \(\Pi\) and minitive \(N\); conditioning is not fixed by those
units and is defined for the example below &
Additivity; complementation \(N(A^c)=1-N(A)\) \\
Cox's theorem (note~5) &
\cref{def:p15-r47-cox} and the boundaries reported after it &
Not a representation but an argument for probabilism from postulates on a
conditional confidence &
No column in the survey's Table~4.1.1; its complement postulate excludes
non-self-dual measures \\
\bottomrule
\end{longtable}

All five rows rest on a choice the survey takes as given: that a single
real number, additive across disjoint events, represents the state. The
following theorem fixes one finite state, described in each framework so
that the static descriptions agree as far as they can, and then applies to
it the learning operations in the table.

\begin{hbtheorem}{thm:p15-w02b-p15-022-alternatives-contrast}{proved=annex:p15-w02b-p15-022-alternatives-contrast}{One state, one piece of evidence, five updates}
Let \(W=\{w_1,w_2,w_3\}\) and \(E=\{w_2,w_3\}\). Let \(\kappa\) be the
negative ranking with world ranks \((0,1,2)\), let
\(\pi=2^{-\kappa}=(1,\tfrac12,\tfrac14)\) on worlds, with
\(\Pi(A)=\max_{w\in A}\pi(w)\), and let \(m\) be the nested mass
assignment \(m(\{w_1\})=\tfrac12\), \(m(\{w_1,w_2\})=\tfrac14\),
\(m(W)=\tfrac14\). Let \(p\) be the pignistic probability of \(m\).
\begin{enumerate}
\item \emph{Static agreement.} \(\operatorname{Pl}_m(\{w\})=\pi(w)\) for
      every \(w\), and \(\operatorname{Bel}_m(A)=N_\kappa(A)\) for every
      \(A\subseteq W\), with \(N_\kappa\) as in \cref{def:p15-r47-ranking}.
      Moreover \(p=(\tfrac{17}{24},\tfrac5{24},\tfrac1{12})\). The events
      believed at threshold zero under \(\kappa\), the events of positive
      necessity, and the events containing the unique minimal world all
      coincide: they are the events containing \(w_1\).
\item \emph{Learning \(E\) with certainty.}
      \begin{enumerate}
      \item Conditioning gives \(p(\cdot\mid E)=(0,\tfrac57,\tfrac27)\).
      \item Ranking conditioning gives \(\kappa(\cdot\mid E)\) world ranks
            \((\infty,0,1)\); the believed events become those containing
            \(w_2\).
      \item Selecting the minimal \(E\)-worlds in the order
            \(w_1\prec w_2\prec w_3\) gives \(\{w_2\}\).
      \item With \(\pi_E(w)=\pi(w)/\Pi(E)\) on \(E\) and \(\pi_E(w_1)=0\),
            \(\pi_E=(0,1,\tfrac12)=2^{-\kappa(\cdot\mid E)}\).
      \item Combining \(m\) with mass one on \(E\) by
            \cref{def:p15-r47-combination} gives \(K=\tfrac12\) and
            \(m_E(\{w_2\})=m_E(E)=\tfrac12\). The singleton
            plausibilities of \(m_E\) are the values of \(\pi_E\);
            the belief of \(\{w_2\}\) under \(m_E\) is \(\tfrac12\); and
            the pignistic probability of \(m_E\) is
            \((0,\tfrac34,\tfrac14)\), not \(p(\cdot\mid E)\).
      \item The probabilities \(q^a=(\tfrac34,0,\tfrac14)\) and
            \(q^b=(\tfrac34,\tfrac14,0)\) both dominate
            \(\operatorname{Bel}_m\) on every event, have
            \(q^a(E)=q^b(E)=\tfrac14\), and give
            \(q^a(\{w_2\}\mid E)=0\), \(q^b(\{w_2\}\mid E)=1\). Conditioning
            the members of \(\{q:q\geq\operatorname{Bel}_m\}\) on \(E\)
            therefore gives \(\{w_2\}\) the range \([0,1]\), whereas the
            combination in~(e) gives \([\tfrac12,1]\).
      \end{enumerate}
\item \emph{Learning the complement afterwards.} After~(a),
      \(p(E^c\mid E)=0\): neither the ratio conditional on \(E^c\) nor a
      Jeffrey update on \(\{E,E^c\}\) is defined. After~(b), \(\neg E\)
      has infinite rank, outside \cref{thm:p15-rat53-update}. Accepting
      \(E\) instead with firmness \(1\) gives world ranks \((1,0,1)\), and
      then accepting \(\{w_1\}\) with firmness \(1\) returns \((0,1,2)\).
      Minimal-world selection for the input \(\{w_1\}\) gives \(\{w_1\}\)
      in every order.
\item \emph{Betting.} Read as two-sided prices for unit bets,
      \(\operatorname{Bel}_m(\{w_1\})+\operatorname{Bel}_m(E)=\tfrac12\).
      Selling both bets at these prices loses \(\tfrac12\) in every world.
\item \emph{Cox's complement postulate.} There is no function \(n\) with
      \(N_\kappa(A^c)=n(N_\kappa(A))\) for all \(A\), and hence none for
      \(\operatorname{Bel}_m\); every probability \(q\) on \(W\) has
      \(q(A^c)=1-q(A)\).
\end{enumerate}
\end{hbtheorem}

\paragraph{Justification of the mapping.}
Item~1 shows that the five descriptions can be made to agree on what is
accepted and on the shape of the uncertainty, so the differences in item~2
come from the update rules, not from different starting information. The
ranking, possibility and minimal-world rules agree: each moves acceptance
to the most plausible world compatible with \(E\), and the ranking and
possibility conditionals are the same function on two scales. This is the
selection mechanism of \cref{thm:p08-agm-representation} applied on the
three-world space, and the comparison made in \cref{sec:p15-rationality-5-3}
between ranking acceptance and faithful-preorder revision. The graded
updates disagree with one another. Conditioning the pignistic probability
gives \(\tfrac57\) to \(w_2\), the pignistic probability of the combined
mass gives \(\tfrac34\), and conditioning every dominating probability
leaves the whole interval \([0,1]\). Projection to one probability does
not commute with any of these updates, so the survey's arguments for
conditioning, which concern a single probability
\cite[p.~244]{HajekStaffel2021Dynamics}, do not transfer to belief
functions through the pignistic projection.

\paragraph{Remaining belief-change comparisons.}
Uncertain learning has three different forms here: Jeffrey updating
reassigns cell probabilities (\cref{def:p15-rat41-jeffrey}); the ranking
update reassigns cell ranks while preserving within-cell conditional ranks
(\cref{def:p15-rat53-soft-update}); a simple support combines a reliability
weight with the existing masses (\cref{thm:p15-r47-combination}). One-step
AGM revision has no graded input. The rules also differ on iteration.
Jeffrey resets do not commute (\cref{thm:p15-rat41-jeffrey}). Item~3 shows
the limitation the survey itself raises about information loss
\cite[p.~246]{HajekStaffel2021Dynamics}: after conditioning, the evidence
has probability one and cannot later be given up by conditioning or by a
Jeffrey update with positive old cells. Finite-firmness ranking updates
are reversible in the example, and AGM revision accepts any consistent
input by its Success and Consistency postulates.

\paragraph{Limitation from the survey's own justifications.}
Item~4 applies the Dutch-book theorem as the survey states it, for
normalized prices that violate the calculus
\cite[p.~240]{HajekStaffel2021Dynamics}: a belief function read as a
schedule of two-sided prices is Dutch-bookable whenever it is not additive.
Belief-function and lower-prevision accounts avoid this conclusion by
reading such numbers as one-sided buying prices
(\cref{sec:p15-r47-lower-previsions}), which the survey's betting account,
with its single buy-or-sell price, excludes by assumption
\cite[p.~240]{HajekStaffel2021Dynamics}. The survey's justification routes
therefore decide against these alternatives only under that interpretation
of credence; they do not decide between the interpretations. Item~5 gives
the corresponding limitation for note~5: Cox's postulates, in the form
recorded in \cref{def:p15-r47-cox}, presuppose that the confidence in a
complement is fixed by the confidence in the event. A necessity measure or
belief function violates that hypothesis outright. Whatever the further
proof boundaries reported there, Cox's theorem bears only on confidence
measures whose complements are so determined, and it cannot by itself rule
out the frameworks of this section.

\section{Reported theorems and their sources}
\label{sec:p15-w02b-p15-022-dependent-sources}

H\'ajek and Staffel report many theorems and one empirical finding without
proving them; each is attributed to a work in their reference list
\cite[pp.~249--251]{HajekStaffel2021Dynamics}. The table fixes the complete
set of works on which those reported claims depend, with the version the
list gives, and records for each claim what is proved in this handbook and
what remains owed to the named work. A finite theorem proved in the
handbook is new mathematics with its own hypotheses; it is not the cited
author's argument, and it does not settle the normative premise the survey
attaches to that argument. The empirical claim of row D21 also points to
two chapters of the same handbook, cited in the text rather than in the
reference list. Those two chapters are held and were examined for this
row; none of the other dependent works is held for this handbook, so none
of their presentations is reconstructed.

\begin{longtable}{@{}p{.28\linewidth}p{.32\linewidth}p{.30\linewidth}@{}}
\caption{Reported claims of the survey and their dependent sources.}
\label{tab:p15-w02b-p15-022-dependent-sources}\\
\toprule Reported claim (survey page) & Dependent work, as listed (list page) & Status here \\
\midrule\endfirsthead
\caption[]{Reported claims of the survey and their dependent sources (continued)}\\
\toprule Reported claim (survey page) & Dependent work, as listed (list page) & Status here \\
\midrule\endhead
\midrule\multicolumn{3}{r}{\emph{continued on next page}}\\
\endfoot
\endlastfoot
D1. Dutch book theorem, stated by Ramsey and proved by de Finetti (p.~240) &
de Finetti, \emph{Foresight: its logical laws, its subjective sources}, in
Kyburg and Smokler (eds.), \emph{Studies in Subjective Probability}, Wiley,
1964, pp.~100--110; French original 1937 (p.~249) &
Finite equivalence proved in \cref{thm:p17-finite-price-coherence} under
its own contract model; de Finetti's proof owed to that work \\
D2. Converse Dutch book theorem (p.~240) &
Kemeny, \emph{Fair bets and inductive probabilities}, J.\ Symbolic Logic
20(3), 1955, 263--273 (p.~250) &
Finite converse in \cref{thm:p17-finite-price-coherence}; Kemeny's
version owed \\
D3. Ramsey's calibration of utility and credence from preferences (p.~241) &
Ramsey, \emph{Truth and probability} (1926), in Mellor (ed.),
\emph{Philosophical Papers}, Cambridge University Press, 1990, pp.~52--94
(p.~250) &
Owed to that work \\
D4. Savage's representation: unique probability, utilities unique up to
positive linear transformation (p.~241) &
Savage, \emph{The Foundations of Statistics}, Wiley, 1954 (p.~250) &
Owed to that work; the finite models of
\cref{thm:p17-finite-bet-representation,thm:p17-finite-vnm} are different
results \\
D5. Jeffrey's axiomatization and Bolker's representation, typically without
a unique probability (p.~241) &
Jeffrey, \emph{The Logic of Decision}, University of Chicago Press, 1965
(p.~249). The list gives no separate entry for Bolker; the survey reports
his theorem as cited by Jeffrey &
Owed to Jeffrey 1965 and to Bolker's own publication, which the survey does
not identify \\
D5a. An agent representable as an expected-utility maximizer may also be
representable nonprobabilistically, or in implausible ways (p.~241) &
Zynda, \emph{Representation theorems and realism about degrees of belief},
Philosophy of Science 67, 2000, 45--69 (p.~251); H\'ajek,
\emph{Arguments for---or against---probabilism?}, British J.\ Philosophy of
Science 59(4), 2008, 793--819 (p.~249) &
Owed to those works \\
D6. Accuracy dominance for Brier scores, first shown by de Finetti
(p.~242) &
de Finetti, \emph{Theory of Probability}, Wiley, 1974, pp.~87--91 as cited
(p.~249) &
Finite Brier dominance and its converse proved in
\cref{thm:p15-rat41-brier}; de Finetti's version owed \\
D7. Extensions of dominance to further scoring rules (p.~242) &
Joyce, \emph{A non-pragmatic vindication of probabilism}, Philosophy of
Science 65(4), 1998, 575--603 (p.~249); Leitgeb and Pettigrew,
\emph{An objective justification of Bayesianism I} and \emph{II},
Philosophy of Science 77(2), 2010, 201--235 and 236--272 (p.~250);
Pettigrew, \emph{Accuracy and the Laws of Chance}, Oxford University
Press, 2016 (p.~250) &
Owed to those works; no result beyond the Brier case is proved here \\
D8. Regularity is impossible for uncountably many possible outcomes
(p.~242) &
H\'ajek, \emph{What conditional probability could not be}, Synthese
137(3), 2003, 273--323 (p.~249) &
Proved for finitely additive probability in
\cref{thm:p15-rat41-positive-atoms}; H\'ajek's presentation owed \\
D8a. A zero credence in a possible event exposes one to a semi-Dutch book
(p.~242) &
Shimony, \emph{Coherence and the axioms of confirmation}, J.\ Symbolic
Logic 20(1), 1955, 1--28 (p.~250) &
The one-bet payoff is computed in \cref{sec:p15-rat41-synchronic};
Shimony's argument owed \\
D9. Accuracy and almost-Dutch-book arguments for the Principal Principle
(p.~243) &
Pettigrew 2016, as in D7; Pettigrew, \emph{An almost-Dutch book argument
for the Principal Principle}, web post of 2 February 2018 (p.~250) &
Owed to those works \\
D10. Dutch book for Reflection (p.~243) &
van Fraassen, \emph{Belief and the will}, J.\ Philosophy 81(5), 1984,
235--256 (p.~251) &
Owed to that work. \Cref{thm:p15-w02b-p15-022-reflection} proves only the
Reflection identity for a self that conditions a fixed exchangeable law,
which is not a Dutch-book argument \\
D11. Obeying Reflection maximizes expected accuracy (p.~243); any coherent
credences expect themselves to be most accurate under a strictly proper
measure (p.~244) &
Easwaran, \emph{Expected accuracy supports conditionalization---and
conglomerability and reflection}, Philosophy of Science 80(1), 2013,
119--142 (p.~249) &
Owed to that work; strict propriety for one event is proved in
\cref{thm:p15-rat41-brier} \\
D12. Minimax accuracy of the indifferent distribution (p.~243) &
Pettigrew 2016, as in D7 &
Finite Brier minimax proved in \cref{thm:p15-rat41-indifference};
Pettigrew's version owed \\
D13. Violating conditionalization exposes one to a diachronic Dutch book
(p.~244) &
Lewis, \emph{Why conditionalize?}, in \emph{Papers in Metaphysics and
Epistemology}, Cambridge University Press, 1999, pp.~403--407 (p.~250) &
Owed to that work; no diachronic Dutch-book theorem is proved here \\
D14. Converse: conditionalizers who commit no other Dutch-bookable offense
are immune (p.~244) &
Skyrms, \emph{Dynamic coherence and probability kinematics}, Philosophy of
Science 54(1), 1987, 1--20 (p.~250) &
Owed to that work \\
D15. A fully specified decision lets the bookie win against an informed
agent (p.~244) &
Skyrms, \emph{A mistake in dynamic coherence arguments?}, Philosophy of
Science 60(2), 1993, 320--328 (p.~250) &
Owed to that work \\
D16. Conditioning maximizes expected accuracy relative to the initial
credences (p.~244) &
Easwaran 2013, as in D11; Greaves and Wallace, \emph{Justifying
conditionalization: conditionalization maximizes expected epistemic
utility}, Mind 115(459), 2006, 607--632 (p.~249); Leitgeb and Pettigrew
2010, part II, as in D7; Pettigrew 2016, as in D7; Schoenfield,
\emph{Conditionalization does not (in general) maximize expected accuracy},
2017 (p.~250) &
Brier planning case on a finite partition proved in
\cref{thm:p15-rat41-brier}; the general arguments, and Schoenfield's
qualification, owed to those works \\
D17. Minimal-distance characterizations of conditioning and Jeffrey
conditioning, and noncommutativity of Jeffrey updating (pp.~244--245) &
Diaconis and Zabell, \emph{Updating subjective probability}, J.\ American
Statistical Association 77(380), 1982, 822--830 (p.~249); Levi,
\emph{Probability kinematics}, British J.\ Philosophy of Science 18(3),
1967, 297--315 (p.~250) &
Finite characterizations proved in \cref{thm:p15-rat41-minimal-update},
with the nonuniqueness of total variation; a noncommuting pair proved in
\cref{thm:p15-rat41-jeffrey}; the cited presentations owed \\
D18. Diachronic Dutch book for Jeffrey conditioning (p.~245) &
Armendt, \emph{Is there a Dutch book argument for probability
kinematics?}, Philosophy of Science 47(4), 1980, 583--588 (p.~249);
Skyrms 1987, as in D14 &
Owed to those works \\
D19. The Judy Benjamin objection to relative-entropy updating (p.~245) &
van Fraassen, \emph{A problem for relative information minimizers in
probability kinematics}, British J.\ Philosophy of Science 32, 1981,
375--379 (p.~251) &
Owed to that work \\
D20. Approximate coherence improves accuracy and Dutch-book loss under
matched measures (p.~247) &
De Bona and Staffel, \emph{Graded incoherence for accuracy-firsters},
Philosophy of Science 84(2), 2017, 189--213, and \emph{Why be
(approximately) coherent?}, Analysis 78(3), 2018, 405--415 (p.~249) &
Squared-Euclidean and Brier case proved in \cref{thm:p15-rat41-graded};
the logarithmic-score comparison owed to those works \\
D21. Empirical finding: people change their credences too little in
response to incoming evidence (p.~246) &
(a) Slovic and Lichtenstein, \emph{Comparison of Bayesian and regression
approaches to the study of information processing in judgment},
Organizational Behavior and Human Performance 6(6), 1971, 649--744
(p.~250). (b) Gl\"ockner, \emph{Prospect Theory}, chapter~8.3 of the same
handbook (cited on p.~246, not in the list). (c) Hertwig and Kozyreva,
\emph{Bounded Rationality: A Vision of Rational Choice in the Real World},
chapter~8.5 of the same handbook (cited on p.~246, not in the list) &
(a) Unheld and owed: the survey gives no task, population, design,
measure or uncertainty, so no finding is stated here. (b) Held and read;
it reports no study of credence revision, so it does not supply the
finding (see below). (c) Held and read; it names Bayesian updating as a
benchmark and reports objections to such benchmarks, but reports no
revision study (see below). The finding stays owed to (a) \\
\bottomrule
\end{longtable}

Rows D1--D21, with D5a and D8a, cover every theorem and finding that the
survey reports as established elsewhere; row D21 lists all three works its
empirical sentence cites. The two held chapters bear on that sentence as
follows.

Gl\"ockner's chapter concerns choice among gambles with stated or sampled
outcomes and probabilities. Its subject is how outcomes and probabilities
are weighted in a choice, measured against expected utility theory, not how
far a credence moves when evidence arrives. The material nearest to belief
revision is its account of decisions from experience, where people sample
outcomes before choosing. There it reports that the direction of the
difference between described and experienced rare events is problem
dependent and does not hold universally \cite[p.~490]{Glockner2021ProspectTheory};
\cref{thm:p17-pt-sampling} and the surrounding discussion treat that
material. None of this is a measurement of posterior credence against a
Bayesian posterior, so the chapter supplies no task, population or
measure for the survey's finding.

Hertwig and Kozyreva describe the heuristics-and-biases program as taking
rational-agent models, with Bayesian probability updating among their
assumptions, as the benchmarks for human judgment
\cite[p.~509]{HertwigKozyreva2021Bounded}. They then report objections to
that program: its norms are narrow coherence standards; some reported
illusions change with the format in which probabilistic information is
presented; and some behavior read as irrational may come from decision
rules adapted to environments other than the laboratory task
\cite[p.~510]{HertwigKozyreva2021Bounded}. These objections are discussed in
\cref{sec:p17-br-heuristics-biases}. They bear on how any underreaction
result should be interpreted, which matches the survey's own remark that
the interpretation of such data is debated \cite[p.~246]{HajekStaffel2021Dynamics}.
The chapter reports no study of the size of credence revisions, however,
so it supplies no finding either. The empirical claim therefore has no
task, population, design, measure, uncertainty estimate or replication
record in any held source. It is recorded here only as the survey's
report, and no normative conclusion is drawn from it.

Its philosophical objections and positions, such as
Christensen's and Levi's criticisms of diachronic Dutch books, are reported
in \cref{sec:p15-rat41-dynamics} as positions, not as claims with proofs.
Each owed entry needs the named work in the listed version, at the pages
carrying the cited argument.

\section{Boundaries of the reverse-limit results}
\label{sec:p15-w02b-p15-022-boundaries}

The reverse-martingale proof of the strong law, the change-of-integrand
theorem and the conditional-projection lemma of
\cref{sec:p15-exchangeability-and-definetti} have the following limiting
cases.

\begin{hbtheorem}{thm:p15-w02b-p15-022-reverse-strong-law-boundary}{proved=annex:p15-w02b-p15-022-reverse-strong-law-boundary}{Where the reverse strong law stops}
Use real coordinates and the fields
\(\mathcal H_n=\sigma(S_n,X_{n+1},X_{n+2},\ldots)\) of
\cref{thm:p15-reverse-strong-law}, with
\(\mathcal H_\infty=\bigcap_n\mathcal H_n\).
\begin{enumerate}
\item If the \(X_j\) are i.i.d. with \(E|X_1|=\infty\), then
      \(P(|X_n|>n\text{ i.o.})=1\), and the event that \(S_n/n\) has a finite
      limit has probability zero.
\item If the law is exchangeable and \(E|X_1|<\infty\), then
      \(\mathcal H_n\subseteq\mathcal E_n\),
      \(S_n/n=E(X_1\mid\mathcal H_n)\), and
      \(S_n/n\to E(X_1\mid\mathcal H_\infty)\) almost surely and in \(L^1\).
\item If \(\xi\) is integrable and not almost surely constant, the law of
      \((\xi,\xi,\ldots)\) is exchangeable, \(S_n/n=X_1\) almost surely, and
      \(\mathcal H_\infty\) contains an event of probability strictly
      between zero and one. The limit is not the constant \(E\xi\).
\end{enumerate}
\end{hbtheorem}

In item~1 the conditional expectation that defines the reverse martingale
does not exist, and the conclusion itself fails; the symmetric example of
\cref{thm:p15-strong-law-application-boundaries} has averages with no limit
at all. Items~2 and~3 show that the reverse-martingale step survives under
exchangeability alone, and that independence enters through the triviality
of the limiting information. That is exactly the gap between the strong law
and the random limit of \cref{thm:p15-w02b-p15-022-predictive}.

\begin{hbtheorem}{thm:p15-w02b-p15-022-backward-changing-boundary}{proved=annex:p15-w02b-p15-022-backward-changing-boundary}{The envelope in the change-of-integrand theorem}
For \cref{thm:p15-backward-changing-variables}:
\begin{enumerate}
\item The integrable envelope cannot be replaced by uniform integrability.
      With the mutually independent Bernoulli variables \(B_n,C_n\) of
      \cref{thm:p15-ui-counterexamples}, \(Z_n=nB_nC_n\), and the constant
      decreasing sequence \(\mathcal G_n=\sigma(B_1,B_2,\ldots)\), the family
      \((Z_n)\) is uniformly integrable and tends to zero almost surely, but
      \(E(Z_n\mid\mathcal G_n)=B_n\) does not converge almost surely. No
      integrable \(W\) dominates every \(|Z_n|\).
\item Without uniform integrability even the means fail: on \((0,1)\) with
      uniform probability, \(Z_n=n\mathbf1_{(0,1/n)}\to0\) everywhere,
      but with trivial \(\mathcal G_n\), \(E(Z_n\mid\mathcal G_n)=1\).
\item If \(Z_n\to Z\) in \(L^1\) and \(\mathcal G_n\) decreases to
      \(\mathcal G_\infty\), then
      \(E(Z_n\mid\mathcal G_n)\to E(Z\mid\mathcal G_\infty)\) in \(L^1\).
      In item~1 the conditional sequence tends to zero in \(L^1\).
\end{enumerate}
\end{hbtheorem}

Item~1 is Durrett's Example~4.6.11 read with a constant, hence decreasing,
sequence of fields \cite[p.~248]{DurrettPTEAuthor2019}. Item~3 is the
decreasing counterpart of the first assertion of
\cref{thm:p15-changing-conditional-limits}.

The reverse-average proof of \cref{thm:p15-reverse-hewitt-savage} rests on
\cref{thm:p15-independent-conditional-projection}: a conditional mean that
has a version measurable with respect to a sigma-field independent of \(V\)
is the constant \(EV\). The next theorem shows what the independence
hypothesis does there, and exactly how much of it the proof uses.

\begin{hbtheorem}{thm:p15-w02b-p15-022-projection-boundary}{proved=annex:p15-w02b-p15-022-projection-boundary}{The independence hypothesis in the conditional-projection lemma}
Let \(V\) be integrable, \(Y=E(V\mid\mathcal G)\), and let \(Y\) have a
version \(Y'\) measurable with respect to a sigma-field \(\mathcal H\).
\begin{enumerate}
\item \emph{Independence cannot be dropped.} If \(V\in L^2\) is not almost
      surely constant and \(\mathcal G=\mathcal H=\sigma(V)\), then \(V\) is
      a version of \(Y\) measurable with respect to \(\mathcal H\),
      \(\mathcal H\) is not independent of \(V\), and
      \(P(Y\neq EV)>0\). For \(P(V=1)=P(V=0)=1/2\):
      \(Y=V\), \(EV=1/2\), and
      \(P(V=1,\,V=1)=1/2\neq1/4=P(V=1)^2\).
\item \emph{The Hewitt--Savage step under exchangeability alone.} Let
      \(P(\xi=1)=P(\xi=0)=1/2\) and \(X_j=\xi\) for every \(j\), on
      \(\{0,1\}^{\mathbb N_+}\). Take \(V=X_1\), \(\mathcal G=\mathcal E\)
      and \(\mathcal H=\sigma(X_2,X_3,\ldots)\). The law is exchangeable;
      \(\bar Y=\limsup_nS_n/n\) is a version of \(E(X_1\mid\mathcal E)\)
      measurable with respect to both \(\mathcal E\) and \(\mathcal H\);
      \(P(X_1=1,\,X_2=1)=1/2\neq1/4\); and \(\bar Y=X_1\neq1/2=EV\) almost
      surely.
\item \emph{What the proof uses.} If
      \(E(V;A)=EV\cdot P(A)\) for every \(A\in\mathcal H\), then
      \(Y=EV\) almost surely; square integrability is not needed. This
      holds when \(\mathcal H\) is independent of \(V\), and it is strictly
      weaker: if \(\varepsilon\) and \(W\) are independent,
      \(P(\varepsilon=1)=P(\varepsilon=-1)=1/2\),
      \(P(W=1)=P(W=2)=1/2\), \(V=\varepsilon W\) and
      \(\mathcal G=\mathcal H=\sigma(W)\), then \(Y=0=EV\) almost surely,
      although \(P(W=1,\ |V|=1)=1/2\neq1/4=P(W=1)P(|V|=1)\).
\end{enumerate}
\end{hbtheorem}

Item~1 is the extreme case: when the conditioning information is the
variable itself, the projection reproduces \(V\), and a field that contains
\(\{V\leq c\}\) cannot be independent of \(V\). Item~2 is the point at which
the reverse-average route of \cref{thm:p15-reverse-hewitt-savage} needs
independent coordinates. There, \(\mathcal H\) is the field of the
coordinates after a deleted initial segment. Under an exchangeable law that is not
i.i.d., the deleted coordinates and the remaining coordinates can share
\(\mathcal E\)-measurable information, so the limit
\(E(X_1\mid\mathcal E)\) is the random directing frequency of
\cref{thm:p15-w02b-p15-022-predictive}, not a constant. Item~3 shows that
all the conclusion needs from independence is that every
\(\mathcal H\)-event is uncorrelated with \(V\); the conclusion then
holds for integrable \(V\), and the given \(\mathcal G\) need not be
independent of \(V\) at all.

\subsection{Mixtures of independent sequences off standard Borel spaces}
\label{sec:p15-w02b-p15-022-presentable}

The conditional-product identity \cref{thm:p15-definetti-factorization}
holds on an arbitrary measurable coordinate space, while the mixture
representation \cref{thm:p15-definetti-kernel} assumes a standard Borel
space. Durrett states the boundary between them. When the coordinates take
values in a nice space, the sequence has a regular conditional distribution
given the exchangeable field and can be represented as a mixture of
i.i.d.\ sequences; Hewitt and Savage (1956) call such a sequence
\emph{presentable}. For an arbitrary measure space he says the
representation is not valid, and he refers to Dubins and Freedman (1979) and
Freedman (1980) for counterexamples \cite[p.~254]{DurrettPTEAuthor2019}.
His bibliography gives Freedman's title as the statement that a mixture of
i.i.d.\ random variables need not admit a regular conditional probability
given the exchangeable field \cite[p.~478]{DurrettPTEAuthor2019}.
His remark on the foundations of statistics, which follows the binary
case on the next page, is discussed after
\cref{thm:p15-w02b-p15-022-predictive}.

Durrett does not define presentability beyond that sentence. The
definition below records the finite-dimensional reading of ``represented
as a mixture of i.i.d.\ sequences'' that is used in this subsection; it is
not Hewitt and Savage's own definition, which is not held.

\begin{hbdefinition}{def:p15-w02b-p15-022-presentable}{Mixture of independent laws on an arbitrary space}
Let \((S,\mathcal S)\) be a nonempty measurable space and let \(P\) be a
law on \(\Omega=S^{\mathbb N_+}\) with the setting of
\cref{def:p15-exchangeable-canonical}. Let \(\mathscr P(S)\) be the set of
probability measures on \(\mathcal S\), with the sigma-field generated by
the evaluations \(e_B(\mu)=\mu(B)\), \(B\in\mathcal S\); for standard Borel
\(S\) this is the space of \cref{def:p15-directing-law}. Call \(P\)
\emph{presentable} if there is a probability \(\Lambda\) on
\(\mathscr P(S)\) such that, for every \(n\geq1\) and
\(B_1,\ldots,B_n\in\mathcal S\),
\[
 P(X_1\in B_1,\ldots,X_n\in B_n)
            =\int_{\mathscr P(S)}\prod_{j=1}^n\mu(B_j)\,\Lambda(d\mu).
\]
A regular conditional law of \(X_1\) given \(\mathcal E\) is a map
\(\omega\mapsto K_\omega\) from \(\Omega\) to \(\mathscr P(S)\) such that
each \(K(B)\) is \(\mathcal E\)-measurable and is a version of
\(P(X_1\in B\mid\mathcal E)\).
\end{hbdefinition}

\begin{hbtheorem}{thm:p15-w02b-p15-022-presentable}{proved=annex:p15-w02b-p15-022-presentable}{Where the regular conditional law enters}
In the setting of \cref{def:p15-w02b-p15-022-presentable}:
\begin{enumerate}
\item Every presentable law is exchangeable, on every measurable space.
\item If \(P\) is exchangeable and \(X_1\) has a regular conditional law
      \(K\) given \(\mathcal E\), then \(K\) is a measurable map into
      \(\mathscr P(S)\), and with \(\Lambda=P\circ K^{-1}\), for every
      \(n\) and \(B_1,\ldots,B_n\in\mathcal S\),
      \[
       P(X_1\in B_1,\ldots,X_n\in B_n\mid\mathcal E)
                    =\prod_{j=1}^nK(B_j)\quad\text{almost surely},
      \]
      and \(P\) is presentable with this \(\Lambda\).
\item If \(S\) is standard Borel, every exchangeable law is presentable.
\item Consequently, an exchangeable law that is not presentable admits no
      regular conditional law of \(X_1\) given \(\mathcal E\).
\end{enumerate}
\end{hbtheorem}

Item~2 locates the nice-space hypothesis exactly: the factorization of
\cref{thm:p15-definetti-factorization} is available on every space, and
the only further ingredient that the mixture needs is a regular version of
the conditional law of one coordinate. Item~4 is the contrapositive, so a
counterexample of the Dubins--Freedman kind must be a space on which that
regular version fails to exist. The converse of item~2 is not asserted:
by the title Durrett reports, Freedman's paper concerns a presentable
sequence without a regular conditional probability given the exchangeable
field. The counterexamples themselves and Hewitt and Savage's definition
are in works not held for this handbook
(Dubins and Freedman, \emph{Exchangeable processes need not be distributed
mixtures of independent and identically distributed random variables},
Z.\ Wahrscheinlichkeitstheorie verw.\ Gebiete 48, 1979, 115--132;
Freedman, Z.\ Wahrscheinlichkeitstheorie verw.\ Gebiete 51, 1980,
239--248; Hewitt and Savage, \emph{Symmetric measures on Cartesian
products}, Trans.\ AMS 80, 1956, 470--501, as Durrett lists them). They
remain owed to those works and are not reconstructed here.

\begin{hbexample}{def:p15-w02b-p15-022-constant-presentable}{A presentable law on an arbitrary space}
Let \(\nu\) be any probability on an arbitrary measurable space
\((S,\mathcal S)\), and let \(P\) be the image of \(\nu\) under
\(x\mapsto(x,x,x,\ldots)\), so every coordinate repeats one random point.
The map \(x\mapsto\delta_x\) into \(\mathscr P(S)\) is measurable, since
\(e_B(\delta_x)=\mathbf1_B(x)\). With \(\Lambda\) the image of \(\nu\)
under that map,
\[
 P(X_1\in B_1,\ldots,X_n\in B_n)=\nu\Bigl(\bigcap_{j=1}^nB_j\Bigr)
     =\int_S\prod_{j=1}^n\delta_x(B_j)\,\nu(dx)
     =\int_{\mathscr P(S)}\prod_{j=1}^n\mu(B_j)\,\Lambda(d\mu).
\]
So this law is presentable on any space, with point-mass components: each
component is the i.i.d.\ law of a constant sequence. For inductive use,
the example is the extreme case of the mixture conclusion: once one
outcome is observed, the mixture predicts that every later outcome
repeats it.
\end{hbexample}

\section{Procedure, relationships and sources}
\label{sec:p15-w02b-p15-022-procedure}

\begin{hbmethod}{def:p15-w02b-p15-022-procedure}{Assess an exchangeability-based inductive inference}
\begin{enumerate}
\item State the instances, their indexing, the outcome description and the
      target instance or limiting quantity.
\item Give the reason for treating the indices as uninformative. Check the
      commitments of \cref{thm:p15-w02b-p15-022-symmetry-consequences}
      against any trend, persistence or periodicity known for the process.
      Where order matters, consider a partial exchangeability or an explicit
      Markov model instead.
\item State whether the population is infinite or extendable. For a finite
      population without replacement, use the finite mixture of
      \cref{thm:p15-carnap-history-finite}.
\item State the mixing law and its support. Report how the prediction
      changes under alternative mixing laws, including other coordinates of
      indifference.
\item State the outcome description, and whether a recoded description
      would also be judged exchangeable.
\item Separate the formal conclusion from the norms invoked: coherence for
      the joint law, conditioning for later credence, and countable
      additivity for limiting-frequency claims.
\item If the state is not a single additive probability, name the
      framework and its update rule from
      \cref{tab:p15-w02b-p15-022-alternatives}, and do not carry the
      survey's betting or accuracy justifications across without the
      interpretation of credence they assume.
\end{enumerate}
\end{hbmethod}

\begin{proofboundary}{def:p15-w02b-p15-022-scope}{Sources and proof status}
Proved here, in the annex: the lemma and the eleven theorems of this
chapter, from \cref{thm:p15-w02b-p15-022-symmetry-consequences} to
\cref{thm:p15-w02b-p15-022-presentable}, from the results of
\cref{sec:p15-exchangeability-and-definetti}, the measure-theoretic
chapters they cite, and the finite belief-change units listed in
\cref{tab:p15-w02b-p15-022-alternatives}.

Reported with locators, not proved: Zabell's account of the permutation
postulate and its predecessors (p.~272), the general multinomial
representation (pp.~275--276), finite extendability (pp.~282--283), and
the Markov-exchangeable representation attributed to Diaconis and Freedman
(p.~289). The comparison with H\'ajek and Staffel applies norms and
objections they state; it does not attribute to them any view about
exchangeability or about the alternative frameworks they name.

Owed to unheld works: every entry of
\cref{tab:p15-w02b-p15-022-dependent-sources} marked as owed, in the
version listed there. Durrett's boundary for arbitrary measurable spaces
is reported with its locator in \cref{sec:p15-w02b-p15-022-presentable};
the Dubins--Freedman (1979) and Freedman (1980) counterexamples and Hewitt
and Savage's 1956 definition of presentability, in the versions Durrett
lists, are unheld and owed to those works. The two held handbook chapters
of row~D21 were read and supply no finding for it.

Not supplied: de Finetti's 1937 presentation, the Diaconis--Freedman Markov
mixture theorem and Savage's foundations discussion cited by Durrett
(second edition, 1972, Section~3.7) are separate original sources, not
held, whose presentations are not reconstructed here.
\end{proofboundary}

\index{exchangeability!justification of induction}
\index{de Finetti representation!inductive use}
\index{Reflection Principle!exchangeable updating}
\index{Indifference Principle!mixing law}
